% =====================================================================
%  3.5 Models
% =====================================================================
\subsection{Models}
\label{sec:models}

All models receive the same cleaned data and the same information for a given weather setting
(\code{src/models.py}).

\paragraph{B0: seasonal naive.} $\hat y_{D,h}=y_{D-7,h}$. It captures the daily and weekly cycle
for free. It is the floor every model must beat, not the main opponent.

\paragraph{B1: per-hour Ridge with degree days.}
Ridge regression~\cite{hoerl1970ridge} with standardised inputs (scaler fitted on training rows
only). Temperature enters through heating and cooling degree days
\begin{equation}
  \HDD(T)=\max(T_h-T,\,0), \qquad \CDD(T)=\max(T-T_c,\,0),
  \label{eq:degree-days}
\end{equation}
computed for each temperature feature, with separate bases $T_h\le T_c$ so that a flat comfort zone
is allowed. Weekday and month are one-hot encoded. We fit 24 separate models, one per local hour,
with a shared $\alpha$ and shared bases. With one shared $\alpha$, this is equivalent to a single Ridge
in which every feature interacts with the hour. It is the simplest way to let the effect of
temperature, weekday and lags differ between, for example, the 08:00 heating peak and the 15:00
cooling peak. We do \emph{not} interpret Ridge coefficients as ``MW per $^\circ$C'': they are shrunk
and shared between correlated features. Explanation uses SHAP on LightGBM instead.

\paragraph{M1: LightGBM.}
Gradient-boosted trees~\cite{ke2017lightgbm} with one \emph{global} model for all 24 hours, hour
being a feature. Each hour has only about 730 training days, and trees need many samples per leaf.
A global model can still split on the hour when the data support it. LightGBM receives raw
temperatures, since trees are invariant to monotone transforms and find thresholds themselves. Both
B1 and M1 minimise squared error, so the comparison is about the function class, not the loss.

\paragraph{Target option.}
The first validation run showed a bias of $-1.3\%$ of mean load in 2017. Our plan, written
before seeing any result, stated that a bias of at least 1\% triggers a second target option. Both
B1 and M1 can therefore predict either the level $y_{D,h}$ or the change since last week
$y_{D,h}-y_{D-7,h}$, and the choice is made on validation like any other hyper-parameter.

\paragraph{Additional variants (added after the first test run).}
Feedback from the lecturer pointed out that B0$\to$B1 changes two things at once, fitting a model
and engineering features, and that B1 (24 models) and M1 (one model) also differ in structure. We
therefore added five variants (\cref{tab:variants}) so that each question is answered by a pair that
differs in one thing only (\code{src/ablation.py}). \textbf{F0} is a minimal input set: hour, weekday,
month, the raw temperature of the setting ($T_t$ or $T_{D-1,08}$) and $y_{D-7,h}$, i.e.\ the
information B0 uses plus the calendar. It has no degree days, holidays, recent lags, morning ratio or
cyclic encodings. One-hot coding of weekday and month is kept for Ridge, because a linear model needs
it to use a category at all. \textbf{F1} is the full set of \cref{tab:features}. The global Ridge is
one model with hour dummies and no hour interactions; the per-hour LightGBM is 24 models sharing one
set of hyper-parameters, each with its own early stopping. All variants are tuned exactly like B1 and
M1 (\cref{sec:tuning}). For the global Ridge the best $\alpha$ fell on the upper edge of the grid
(one model on 24 times more rows needs a larger penalty), so the grid was extended by half decades,
as our rule prescribes; the chosen value lies inside the extended grid.

\begin{table}[H]
  \centering
  \small
  \caption{Model grid. Each variant exists in both weather settings.}
  \label{tab:variants}
  \begin{tabular}{@{}l l l l@{}}
    \toprule
    Model & Features & Structure & Status \\
    \midrule
    Ridge & F1 & 24 per-hour models & original B1 \\
    LightGBM & F1 & one global model & original M1 \\
    Ridge & F0 & 24 per-hour models & added \\
    Ridge & F0 & one global model & added \\
    Ridge & F1 & one global model & added \\
    LightGBM & F0 & one global model & added \\
    LightGBM & F1 & 24 per-hour models & added \\
    \bottomrule
  \end{tabular}
\end{table}
